% id: 21-18
% book: 베이직쎈 확률과 통계 · 02 중복조합과 이항정리
% section: 개념 04 중복조합
% passage: 같은 종류의 음료수 5개를 3명의 학생에게 나누어 주려고 한다. 다음을 구하시오.
% figure: none
% answer: 
% answer_source: 
% variant_level: 0 (원본 전사)
% 이 파일은 build.mjs 가 items.json 에서 만든다. 고칠 때는 items.json 을 고친다.
\smash{\raisebox{1.5mm}{\dmkichul{베이직쎈 확률과 통계 21쪽 18번}}}
음료수를 학생들이 \underline{적어도 한 개씩 받도록} 나누어 주는 경우의 수
\dmhint{음료수를 학생들이 적어도 한 개씩 받으려면 먼저 음료수를 3명의 학생에게 각각 1개씩 나누어 주고, 나머지 \blank개의 음료수를 3명의 학생에게 나누어 주면 된다.\\ 따라서 구하는 경우의 수는 서로 다른 \blank개에서 \blank개를 택하는 중복조합의 수와 같으므로\\ ${}_{\text{\scriptsize\setlength{\fboxsep}{1.5pt}\framebox[1.5em]{\rule{0pt}{1.6ex}}}}\mathrm{H}_{\text{\scriptsize\setlength{\fboxsep}{1.5pt}\framebox[1.5em]{\rule{0pt}{1.6ex}}}}={}_{\text{\scriptsize\setlength{\fboxsep}{1.5pt}\framebox[1.5em]{\rule{0pt}{1.6ex}}}}\mathrm{C}_2=\blank$}
